% GENERATED FILE. Edit canonical lessons, then run npm run generate:books.
% canonical-source-hash: fnv1a64:549fa6eaa266f183
% canonical-lessons: IT-C360-dirigente, IT-C360-reparto, IT-C360-mansione, IT-C360-proroga, IT-C360-consegna

\chapter{In ufficio}
\label{ch:it-in-ufficio}
\hlchaptermodality{360}

\begin{chapteropening}
\textbf{By the end of this chapter:} \emph{I can describe my place at work: the manager, my department and my duties, a delivery date and an extension.}

Complete the last lesson of chapter 360: I can describe my place at work: the manager, my department and my duties, a delivery date and an extension.
\end{chapteropening}

\section[il dirigente]{il dirigente — the manager, the executive}

\label{lesson:IT-C360-dirigente}



\begin{quote}
Before the new one: say the Italian for the shift, then the Italian for the overtime.
\end{quote}

\subsection*{You'll want to know: il dirigente}
\textbf{il dirigente} — \textquotedblleft{}the manager, the executive\textquotedblright{}.

\begin{quote}
Il dirigente vuole parlare con tutto il gruppo.
\end{quote}

\begin{quote}
\emph{The manager wants to talk to the whole team.}
\end{quote}

Higher than \emph{il capo} of a small team. \emph{Un dirigente scolastico} is a head teacher.

\begin{cousinweb}
\emph{dirigente} is the present participle of \emph{dirigere} \textquotedblleft{}to direct\textquotedblright{}, from Latin \emph{dīrigere} \textquotedblleft{}to set straight, to guide\textquotedblright{} (Lewis \& Short), built on \emph{regere} \textquotedblleft{}to rule, to steer\textquotedblright{}.
\end{cousinweb}

\begin{grammarlens}[title={What you've built}]
One more word for telling, explaining and describing at B1.
\end{grammarlens}

\subsection*{Your turn}
\begin{itemize}
\raggedright
  \item \emph{Say it:} \emph{il dirigente}
  \item \emph{Say it:} \emph{il dirigente}, once more
  \item \emph{Say it:} a sentence of your own with \emph{il dirigente}
  \item \emph{Recall:} say the Italian for the shift, then the Italian for the overtime, then say \emph{il dirigente} again
\end{itemize}

\begin{tcolorbox}[breakable,colback=teal!4,colframe=teal!35!black,title={Before you move on}]
What does il dirigente mean? (\textquotedblleft{}The manager, the executive\textquotedblright{}.) Say it once more.
\end{tcolorbox}
\section[il reparto]{il reparto — the department, the section}

\label{lesson:IT-C360-reparto}



\begin{quote}
Before the new one: say the Italian for the overtime, then the Italian for the manager.
\end{quote}

\subsection*{You'll want to know: il reparto}
\textbf{il reparto} — \textquotedblleft{}the department, the section\textquotedblright{}.

\begin{quote}
Lavoro nel reparto vendite.
\end{quote}

\begin{quote}
\emph{I work in the sales department.}
\end{quote}

In a hospital, \emph{il reparto} is the ward; in a shop, the department.

\begin{cousinweb}
\emph{reparto} is built on \emph{ripartire} \textquotedblleft{}to share out\textquotedblright{}, from \emph{ri-} and \emph{partire} \textquotedblleft{}to divide\textquotedblright{}, inherited from Latin \emph{partīrī} \textquotedblleft{}to divide\textquotedblright{}, from \emph{pars, partis} \textquotedblleft{}a part\textquotedblright{} (Lewis \& Short).
\end{cousinweb}

\begin{grammarlens}[title={What you've built}]
One more word for telling, explaining and describing at B1.
\end{grammarlens}

\subsection*{Your turn}
\begin{itemize}
\raggedright
  \item \emph{Say it:} \emph{il reparto}
  \item \emph{Say it:} \emph{il reparto}, once more
  \item \emph{Say it:} a sentence of your own with \emph{il reparto}
  \item \emph{Recall:} say the Italian for the overtime, then the Italian for the manager, then say \emph{il reparto} again
\end{itemize}

\begin{tcolorbox}[breakable,colback=teal!4,colframe=teal!35!black,title={Before you move on}]
What does il reparto mean? (\textquotedblleft{}The department, the section\textquotedblright{}.) Say it once more.
\end{tcolorbox}
\section[la mansione]{la mansione — the duty, the task (in a job)}

\label{lesson:IT-C360-mansione}



\begin{quote}
Before the new one: say the Italian for the manager, then the Italian for the department.
\end{quote}

\subsection*{You'll want to know: la mansione}
\textbf{la mansione} — \textquotedblleft{}the duty, the task (in a job)\textquotedblright{}.

\begin{quote}
Le mie mansioni sono rispondere al telefono e organizzare le riunioni.
\end{quote}

\begin{quote}
\emph{My duties are answering the phone and organising meetings.}
\end{quote}

Mostly plural in job descriptions: \emph{le mansioni}.

\begin{cousinweb}
\emph{mansione} is usually traced to Latin \emph{mānsiō} \textquotedblleft{}a stay, a lodging\textquotedblright{} (Lewis \& Short), from \emph{manēre} \textquotedblleft{}to remain\textquotedblright{}. How it came to mean a duty, perhaps through French, is not settled.
\end{cousinweb}

\begin{grammarlens}[title={What you've built}]
One more word for telling, explaining and describing at B1.
\end{grammarlens}

\subsection*{Your turn}
\begin{itemize}
\raggedright
  \item \emph{Say it:} \emph{la mansione}
  \item \emph{Say it:} \emph{la mansione}, once more
  \item \emph{Say it:} a sentence of your own with \emph{la mansione}
  \item \emph{Recall:} say the Italian for the manager, then the Italian for the department, then say \emph{la mansione} again
\end{itemize}

\begin{tcolorbox}[breakable,colback=teal!4,colframe=teal!35!black,title={Before you move on}]
What does la mansione mean? (\textquotedblleft{}The duty, the task (in a job)\textquotedblright{}.) Say it once more.
\end{tcolorbox}
\section[la proroga]{la proroga — the extension (of a deadline)}

\label{lesson:IT-C360-proroga}



\begin{quote}
Before the new one: say the Italian for the department, then the Italian for the duty.
\end{quote}

\subsection*{You'll want to know: la proroga}
\textbf{la proroga} — \textquotedblleft{}the extension (of a deadline)\textquotedblright{}.

\begin{quote}
Ho chiesto una proroga di due settimane.
\end{quote}

\begin{quote}
\emph{I asked for a two-week extension.}
\end{quote}

Stress on the first syllable: \emph{pròroga}. The verb is \emph{prorogare}.

\begin{cousinweb}
\emph{proroga} is built on \emph{prorogare} \textquotedblleft{}to extend\textquotedblright{}, a learned borrowing from Latin \emph{prōrogāre} \textquotedblleft{}to prolong, to put off\textquotedblright{} (Lewis \& Short), built from \emph{prō} and \emph{rogāre} \textquotedblleft{}to ask\textquotedblright{}.
\end{cousinweb}

\begin{grammarlens}[title={What you've built}]
One more word for telling, explaining and describing at B1.
\end{grammarlens}

\subsection*{Your turn}
\begin{itemize}
\raggedright
  \item \emph{Say it:} \emph{la proroga}
  \item \emph{Say it:} \emph{la proroga}, once more
  \item \emph{Say it:} a sentence of your own with \emph{la proroga}
  \item \emph{Recall:} say the Italian for the department, then the Italian for the duty, then say \emph{la proroga} again
\end{itemize}

\begin{tcolorbox}[breakable,colback=teal!4,colframe=teal!35!black,title={Before you move on}]
What does la proroga mean? (\textquotedblleft{}The extension (of a deadline)\textquotedblright{}.) Say it once more.
\end{tcolorbox}
\section[la consegna]{la consegna — the delivery, the handing in}

\label{lesson:IT-C360-consegna}



\begin{quote}
Before the new one: say the Italian for the duty, then the Italian for the extension.
\end{quote}

\subsection*{You'll want to know: la consegna}
\textbf{la consegna} — \textquotedblleft{}the delivery, the handing in\textquotedblright{}.

\begin{quote}
La consegna del progetto è fissata per lunedì.
\end{quote}

\begin{quote}
\emph{The handing in of the project is set for Monday.}
\end{quote}

\emph{La data di consegna} — the delivery date, the deadline. \emph{Consegna a domicilio} — home delivery.

\begin{cousinweb}
\emph{consegna} is built on \emph{consegnare} \textquotedblleft{}to hand over\textquotedblright{}, inherited from Latin \emph{cōnsignāre} \textquotedblleft{}to seal, to put on record\textquotedblright{} (Lewis \& Short), built on \emph{signum} \textquotedblleft{}a mark, a seal\textquotedblright{}.
\end{cousinweb}

\begin{grammarlens}[title={What you've built}]
One more word for telling, explaining and describing at B1.
\end{grammarlens}

\subsection*{Your turn}
\begin{itemize}
\raggedright
  \item \emph{Say it:} \emph{la consegna}
  \item \emph{Say it:} \emph{la consegna}, once more
  \item \emph{Say it:} a sentence of your own with \emph{la consegna}
  \item \emph{Recall:} say the Italian for the duty, then the Italian for the extension, then say \emph{la consegna} again
\end{itemize}

\begin{tcolorbox}[breakable,colback=teal!4,colframe=teal!35!black,title={Before you move on}]
What does la consegna mean? (\textquotedblleft{}The delivery, the handing in\textquotedblright{}.) Say it once more.
\end{tcolorbox}
